\originalsection{不等振幅、斜入射和光程的精确几何}\label{tb:young-geometry}
两同频同偏振相干波的振幅$E_{10},E_{20}$不同, 总场平方的时间平均给
\[I=I_1+I_2+2\sqrt{I_1I_2}\cos\phi.\]
所以$I_{\max}=(\sqrt{I_1}+\sqrt{I_2})^2$, $I_{\min}=(\sqrt{I_1}-\sqrt{I_2})^2$. 不等强光的暗纹未必全黑. 两线偏振夹角$\alpha$时交叉项再乘$\cos\alpha$, 垂直偏振不会给这类强度干涉条纹. 若随机相位在探测时间内遍历, 平均余弦为零, 强度直接相加.

两缝横坐标$\pm d/2$, 屏点$(L,y)$, 距离平方为$r_\pm^2=L^2+(y\mp d/2)^2$. 相减再因式分解得$r_- -r_+=2yd/(r_-+r_+)\simeq d\sin\theta$, 其中$\sin\theta=y/\sqrt{L^2+y^2}$, 要求$L$远大于孔径尺度. 进一步$|y|\ll L$才变为$dy/L$. 斜入射光先在两缝带入相位差, 原14.11.3约定的光程差为$d(\sin\theta_1-\sin\theta_2)$, 亮纹满足它等于$m\lambda$. 改变哪束先相减会把$m$反号, 亮纹的物理集合不变.

\begin{example}
\textbf{原14.1: 屏点究竟亮还是暗.} $d=0.150\,\mathrm{mm},L=120\,\mathrm{cm},\lambda=833\,\mathrm{nm},y=2.00\,\mathrm{cm}$. 求光程差、它相对波长的比例和干涉状态.
\end{example}
\begin{solution}
小角近似$\delta=dy/L=2.50\times10^{-6}\,\mathrm m$, $\delta/\lambda=3.0012$. 题给精度下约为3个波长, 属第3级亮纹. 以$\sin\theta=y/\sqrt{L^2+y^2}$更精确计算给$2.49965\,\mu\mathrm m$, 仍在题给精度内接近3. 同一组有限有效数字不能用超出其精度的微小偏离宣称严格数学意义的精确极大.
\end{solution}

\originalsection{三缝、N缝与单缝是同一相量求和}\label{tb:n-slit}
\begin{example}
\textbf{原14.2: 三条等幅窄缝.} 三缝相邻中心距$d$, 相邻相位差$\phi=2\pi d\sin\theta/\lambda$, 每缝场幅$E_0$. 求强度及主次峰比.
\end{example}
\begin{solution}
首尾场相加给$2E_0\cos\phi\sin(\omega t+\phi)$, 再加中间场, 总场为$E_0(1+2\cos\phi)\sin(\omega t+\phi)$. 因此
\[\frac I{I_{\max}}=\frac{(1+2\cos\phi)^2}9.\]
主峰在$\phi=2\pi m$, 场幅$3E_0$, 次峰在$(2m+1)\pi$, 场幅大小$E_0$, 强度比主:次为$9:1$. 零点在$\cos\phi=-1/2$, 每对主峰间有两个零点和一个次峰. 三条波的强度不能直接按三倍相加, 那会丢掉相干交叉项.
\end{solution}

$N$缝等幅场的相量和是$E_0\sum_{j=0}^{N-1}e^{ij\phi}$. 对$e^{i\phi}\ne1$, 几何级数给
\[\sum_{j=0}^{N-1}e^{ij\phi}=e^{i(N-1)\phi/2}\frac{\sin(N\phi/2)}{\sin(\phi/2)},\quad
\frac I{I_{\max}}=\frac1{N^2}\left[\frac{\sin(N\phi/2)}{\sin(\phi/2)}\right]^2.\]
$\phi=2\pi m$时原始和为$N$, 由极限取主峰为1. 主峰两侧最近零点相差相位$2\pi/N$, 近非掠射主峰的半宽量级$\Delta\theta\simeq\lambda/(Nd\cos\theta_m)$. 光栅缝越多峰越窄; 可见阶数还必须$|m\lambda/d|\le1$. 分辨不同波长要比较它们的角分离与峰宽, 不能只因都有主峰就必能分开.

把单缝宽$a$分成$N$条小带, 相邻相位$\Delta\phi=ka\sin\theta/N$, 每带场幅为中央总幅的$1/N$. 总相位跨度$2\beta=ka\sin\theta$, 求和后归一化幅度为
\[\frac1N\frac{\sin\beta}{\sin(\beta/N)}\xrightarrow{N\to\infty}\frac{\sin\beta}{\beta}.\]
这就是前节连续孔径积分的离散极限. 原14.10相量多边形在有限$N$时应按弦长求和, 只有连续极限才把小弦近似成弧长; 本稿用精确级数后取极限, 不把有限多边形的弦弧相等当精确关系.

\originalsection{单缝暗纹的配对图像与旁瓣极值}\label{tb:single-slit-extrema}
当$a\sin\theta=\lambda$时, 上半缝一点与下半缝隔$a/2$的对应点路程差$\lambda/2$, 两小元同幅相消. 一对对配起来得到第一个暗纹. 第$m$个暗纹可分成$2|m|$份, 每相邻对应份相差半波, 得$a\sin\theta=m\lambda$($m\ne0$). $m=0$时无半波相消, 而是中央最亮. 第一对暗纹间中央区小角宽度$\Delta y\simeq2\lambda L/a$, 窄缝使其变宽.

\begin{example}
\textbf{原14.3: 用第一暗纹定位屏.} $\lambda=600\,\mathrm{nm}$, 单缝宽$a=0.800\,\mathrm{mm}$, 第一暗纹离中央$y_1=1.00\,\mathrm{mm}$. 求屏距和中央亮区宽度.
\end{example}
\begin{solution}
$\sin\theta_1=\lambda/a=7.50\times10^{-4}$. 因$y_1=L\tan\theta_1$, $L=y_1\sqrt{1-(\lambda/a)^2}/(\lambda/a)\simeq1.33\,\mathrm m$. 中央两第一暗纹之间宽度$2y_1=2.00\,\mathrm{mm}$. 小角公式与精确几何在此给相同的三位有效数字.
\end{solution}

旁瓣的\emph{精确}局部极大不能都取在相邻暗纹正中. 对$I/I(0)=(\sin\beta/\beta)^2$求导, 除去已知暗纹和中央极限, 得$\beta\cos\beta-\sin\beta=0$, 即$\tan\beta=\beta$. 第一正旁瓣根约$4.4934$, 第二约$7.7253$. 半整数$\beta\simeq(m+1/2)\pi$只是近似, 越高阶相对误差越小.

\begin{example}
\textbf{原14.11.5: 第二旁瓣求波长.} 缝宽$a=0.800\,\mathrm{mm}$, 屏距$L=0.800\,\mathrm m$, 第二旁亮峰距中央$y=1.60\,\mathrm{mm}$. 求波长并说明原教材近似.
\end{example}
\begin{solution}
原件把第二旁峰放在第2、第3暗纹中点, 以$y\simeq(5/2)\lambda L/a$得$\lambda\simeq640\,\mathrm{nm}$. 用实际极大条件$\beta_2=7.7253$及$\sin\theta=y/\sqrt{L^2+y^2}$, 得$\lambda=\pi a\sin\theta/\beta_2\simeq651\,\mathrm{nm}$. 原640是所声明中点近似的答案, 不是精确旁瓣位置. 本稿数值根另实算留存, 数值核验不代替求导得到的普遍条件.
\end{solution}

\originalsection{双缝的级数计数与完整例题}\label{tb:interference-examples}
\begin{example}
\textbf{原14.11.1、14.11.2、14.11.4.} (a)$d=0.320\,\mathrm{mm},\lambda=500\,\mathrm{nm}$, 数$-45^\circ<\theta<45^\circ$中的理想窄双缝亮纹. (b)$d=0.100\,\mathrm{mm},L=1.00\,\mathrm m,\lambda=500\,\mathrm{nm}$, 求$\theta=0.800^\circ$和$y=4.00\,\mathrm{mm}$的相位差, 以及相位$1/3\,\mathrm{rad}$和光程$\lambda/4$对应的角度. (c)同幅双缝某点为最大强度的60.0\%, 求最小非负相位差及$\lambda=500\,\mathrm{nm}$时的光程差.
\end{example}
\begin{solution}
(a)$|m|<d\sin45^\circ/\lambda=452.548$, 所以$m=-452,\ldots,452$, 共905个. 区间严格端点影响整数计数, 本题端点不是整数故不额外删纹.

(b)以$\phi=2\pi d\sin\theta/\lambda$算, $0.800^\circ$给$17.5\,\mathrm{rad}$. 屏点$y=4.00\,\mathrm{mm}$小角给$\phi\simeq2\pi dy/(\lambda L)=5.03\,\mathrm{rad}$. 反解$\phi=1/3$有$\theta=\arcsin[\lambda/(6\pi d)]=0.0152^\circ$. 光程$\lambda/4$有$\theta=\arcsin[\lambda/(4d)]=0.0716^\circ$. 相位可以超过$2\pi$, 讨论相同强度时才按整周等价.

(c)$\cos^2(\phi/2)=0.600$, 最小非负值$\phi=2\arccos\sqrt{0.600}=1.3694\,\mathrm{rad}$. $\delta=\lambda\phi/(2\pi)=109\,\mathrm{nm}$. 所有更大解还有整周和平镜对称, 题目限定最小值避免多解混淆.
\end{solution}

若中央单缝包络位于第一暗纹内$|\sin\theta|<\lambda/a$, 名义双缝增强级数满足$|m|<d/a$, 还受可见角$|m\lambda/d|\le1$限制. $d/a=n$整数且无额外角截断时有$2n-1$个中央名义级数, 位于边界$m=\pm n$的增强级正好缺级. 非整数时不能照套同一整数公式. 观察到7个中央名义级并不唯一推出$d/a=4$, 例如$3<d/a\le4$在相同无角截断条件下都可有7级.

\begin{example}
\textbf{原14.11.6: 有限宽双缝的缺级.} $\lambda=500\,\mathrm{nm}$, 缝宽$a=0.700\,\mu\mathrm m$, 中心距$d=2.80\,\mu\mathrm m$, 半圆柱屏收集$-90^\circ\le\theta\le90^\circ$. 求名义增强方向、实际非零增强级数及各方向相对中央强度.
\end{example}
\begin{solution}
增强方向$\sin\theta_m=m\lambda/d$, 允许$|m|\le5$, 共11方向: $0$, $\pm10.29^\circ$, $\pm20.92^\circ$, $\pm32.39^\circ$, $\pm45.58^\circ$, $\pm63.23^\circ$. 因$d/a=4$, $m=\pm4$位于单缝第一暗纹而消失, 所以原题名义非零级数为9. 各名义方向强度
\[\frac{I(\theta_m)}{I(0)}=\left[\frac{\sin(\pi m/4)}{\pi m/4}\right]^2.\]
$m=0,1,2,3,4,5$依次为$1,0.811,0.405,0.0901,0,0.0324$, 负级相同. \emph{这些是干涉因子增强方向处的强度}, 完整包络乘积的局部峰一般有小偏移, 包络边缘还可能出现别的弱极值. “9亮纹”是教材按干涉级数扣缺级的计数, 并非所有精确局部极大点的严格总数.
\end{solution}

\begin{center}
\begin{tikzpicture}
\begin{axis}[width=.9\linewidth,height=5cm,xlabel={$u=d\sin\theta/\lambda$},ylabel={$I/I(0)$},xmin=-5.6,xmax=5.6,ymin=0,ymax=1.08,samples=601]
\addplot[thick,structurecolor,domain=-5.6:5.6] {abs(x)<.0001 ? 1 : (sin(deg(pi*x/4))/(pi*x/4))^2*(cos(deg(pi*x)))^2};
\addplot[dashed,domain=-5.6:5.6] {abs(x)<.0001 ? 1 : (sin(deg(pi*x/4))/(pi*x/4))^2};
\end{axis}
\end{tikzpicture}
\end{center}
实线是双缝总强度, 虚线是单缝包络, 横轴是无量纲观察方向而非实际长度. 在$u=\pm4$包络零点造成缺级, 图也显示包络内真实峰会偏离整数方向. 这让计数约定与实际函数有所区别.

白光含多个波长. 对等光程的中央方向, 每个波长都同相增强, 中央接近白色; 离中心后不同$\lambda$的亮纹角度不同, 侧面出现分色, 再远处不同级数重叠使对比降低. 稳定条纹仍需要同一波长分量在两条路径间足够相干, 不能把两个独立白光灯的强度直接画成稳定双缝图.

改变单缝宽度时还要说明保持什么量. 固定入射场幅和照度、相同观测几何时, 中央场是孔径小元幅度积分而正比$a$, 中央强度正比$a^2$; 中央角宽同时约反比$a$. 完整通过缝的总功率正比$a$, 峰高与总能量并非同一量. 前面的$I/I(0)$图归一化后都在中央为1, 没有显示改变$a$后绝对峰高的变化.
